\ifdefined\gkver\else\def\gkver{student}\fi
\documentclass[\gkver]{gaokaozhenti}
\begin{document}

\chapter{2026年北京}

\begin{enumerate}
\item
%% number: 1
%% paperName: 2026年北京高考第 1 题 3 分
%% typeId: 1
%% type: 单选题
%% chapter: 光学
%% point: 光的折射
%% method: 无
%% score: 3
%% degree: 950
%% duplicateId: 0
%% body:
如图所示，插入水中的筷子看上去产生了侧移。这种现象属于 \xzanswer{B}
\onepicture[w=2.4cm]{figs/01.jpg}
\fourchoices[answer=B]
{光的全反射}
{光的折射}
{光的干涉}
{光的衍射}

%% answer: B
%% memo:
\memoanswer{
插入水中的筷子，其水下部分反射的光线传播时，传播方向会发生偏折，是光的折射现象。

人眼逆着折射光线看去，会认为光线沿直线传播，因此感知到的筷子虚像位置和实际位置发生偏移，造成筷子看上去侧移的效果。

故选 B。
}
\item
%% number: 2
%% paperName: 2026年北京高考第 2 题 3 分
%% typeId: 1
%% type: 单选题
%% chapter: 原子和原子核
%% point: 人工核反应
%% method: 无
%% score: 3
%% degree: 950
%% duplicateId: 0
%% body:
氟 18 是正电子发射断层成像（PET）中常用的放射性核素，可通过 $\ce{^{1}_{1}H}+\ce{^{18}_{8}O}\to\ce{^{18}_{9}F}+$（　）核反应制备，括号内的粒子为 \xzanswer{D}
\fourchoices[answer=D]
{$\ce{^{4}_{2}He}$}
{$\ce{^{0}_{-1}e}$}
{$\ce{^{0}_{1}e}$}
{$\ce{^{1}_{0}n}$}

%% answer: D
%% memo:
\memoanswer{
根据题意，设括号内未知粒子为 $^{A}_{Z}\mathrm{X}$，根据质量数守恒和电荷数守恒可知，$A=1$，$Z=0$。则括号内未知粒子为中子 $\ce{^{1}_{0}n}$。

故选 D。
}
\item
%% number: 3
%% paperName: 2026年北京高考第 3 题 3 分
%% typeId: 1
%% type: 单选题
%% chapter: 气体性质
%% point: 玻意耳定律
%% method: 无
%% score: 3
%% degree: 850
%% duplicateId: 0
%% body:
一定质量的理想气体在等温膨胀过程中 \xzanswer{A}
\fourchoices[answer=A]
{气体压强减小}
{外界对气体做正功}
{气体分子平均动能增大}
{气体与外界无热量交换}

%% answer: A
%% memo:
\memoanswer{
A．一定质量的理想气体等温过程满足玻意耳定律 $pV=C$（$C$ 为常量），膨胀时体积 $V$ 增大，因此压强 $p$ 减小，故 A 正确；

B．气体膨胀时体积增大，是气体对外界做功，外界对气体做负功，故 B 错误；

C．理想气体分子平均动能仅由温度决定，等温过程温度不变，分子平均动能不变，故 C 错误；

D．理想气体内能仅与温度有关，等温过程内能变化 $\Delta U=0$；根据热力学第一定律 $\Delta U=Q+W$（$W$ 为外界对气体做的功），本题中 $W<0$，可得 $Q=-W>0$，即气体从外界吸收热量，存在热量交换，故 D 错误。

故选 A。
}
\item
%% number: 4
%% paperName: 2026年北京高考第 4 题 3 分
%% typeId: 1
%% type: 单选题
%% chapter: 运动和力
%% point: 牛顿定律的应用
%% method: 无
%% score: 3
%% degree: 750
%% duplicateId: 0
%% body:
质量为 $2.0\times10^{3}\Ukg$ 的汽车在水平直道上进行刹车性能测试。汽车速度为 $30\Ums$ 时开始刹车直至停止，刹车过程中所受合力为 $2.0\times10^{4}\UN$，则刹车距离为 \xzanswer{C}
\fourchoices[answer=C]
{$135\Um$}
{$90\Um$}
{$45\Um$}
{$15\Um$}

%% answer: C
%% memo:
\memoanswer{
根据牛顿第二定律 $F=ma$，代入数据得 $a=\frac{F}{m}=\frac{2.0\times10^{4}\UN}{2.0\times10^{3}\Ukg}=10\Umsq$，加速度方向与运动方向相反。

刹车过程是末速度为 $0$ 的匀减速直线运动，由匀变速直线运动速度位移公式 $0-v_{0}^{2}=-2ax$，

整理得 $x=\frac{v_{0}^{2}}{2a}$。

代入 $v_{0}=30\Ums$、$a=10\Umsq$，计算得 $x=\frac{30^{2}}{2\times10}\Um=45\Um$。

故选 C。
}
\item
%% number: 5
%% paperName: 2026年北京高考第 5 题 3 分
%% typeId: 1
%% type: 单选题
%% chapter: 机械波
%% point: 波的图象
%% method: 无
%% score: 3
%% degree: 750
%% duplicateId: 0
%% body:
一列简谐横波在 $t=0$ 时的波形图如图所示，此时 $x=3\Um$ 处的质点 P 位于平衡位置且沿 $y$ 轴正方向运动，经过 $1\Us$ 第一次回到平衡位置。下列说法正确的是 \xzanswer{D}
\onepicture[w=6.0cm]{\includesvg{figs/05}}
\fourchoices[answer=D]
{波沿 $x$ 轴正方向传播}
{波的传播速度为 $4\Ums$}
{一个周期内质点 P 运动的路程为 $4\Um$}
{$t=0.5\Us$ 时，质点 P 的加速度沿 $y$ 轴负方向}

%% answer: D
%% memo:
\memoanswer{
A．此时 $x=3\Um$ 处的质点 P 位于平衡位置且沿 $y$ 轴正方向运动，根据“上下坡法”可知，波沿 $x$ 轴负方向传播，A 错误；

B．从波形图可直接读出波长 $\lambda=4\Um$。

$t=0$ 时 P 在平衡位置且沿 $y$ 轴正方向运动，它离开平衡位置后第一次回到平衡位置的过程是：向上运动到波峰、再向下返回平衡位置，该过程对应半个周期，即 $\frac{T}{2}=1\Us$，得周期 $T=2\Us$。根据波速公式 $v=\frac{\lambda}{T}=\frac{4}{2}\Ums=2\Ums$，B 错误；

C．从图中读出振幅 $A=10\Ucm=0.1\Um$，简谐运动中质点一个周期内的路程为 $4A=0.4\Um$，C 错误；

D．$t=0.5\Us=\frac{T}{4}$，$t=0$ 时 P 从平衡位置沿 $y$ 轴正方向运动，经过 $\frac{T}{4}$ 恰好到达正向最大位移处。简谐运动的加速度方向与位移方向相反，指向平衡位置，因此此时加速度沿 $y$ 轴负方向，D 正确。

故选 D。
}
\item
%% number: 6
%% paperName: 2026年北京高考第 6 题 3 分
%% typeId: 1
%% type: 单选题
%% chapter: 运动和力
%% point: 动量和冲量
%% method: 无
%% score: 3
%% degree: 750
%% duplicateId: 0
%% body:
如图所示，一物块在固定斜面上匀速下滑，不计空气阻力。物块在下滑过程中 \xzanswer{A}
\onepicture[w=4.6cm]{\includesvg{figs/06}}
\fourchoices[answer=A]
{所受合力的冲量为零}
{所受支持力的冲量为零}
{所受重力与摩擦力的合力做正功}
{机械能守恒}

%% answer: A
%% memo:
\memoanswer{
A．合外力冲量等于物块动量变化，匀速运动速度不变、动量变化量为 $0$，因此物块所受合力的冲量为零，A 正确；

B．支持力大小不为零，物块下滑过程的作用时间不为零，根据冲量定义 $I_{N}=Nt$，支持力的冲量不为零，B 错误；

C．物块三力平衡，重力、支持力、摩擦力的总合力为 $0$，因此重力与摩擦力的合力与支持力等大反向，方向垂直斜面向下，而物块位移沿斜面向下，该合力与位移方向垂直，做功为 $0$，C 错误；

D．物块在下滑过程中，动能不变，重力势能减小，则机械能减小，D 错误。

故选 A。
}
\item
%% number: 7
%% paperName: 2026年北京高考第 7 题 3 分
%% typeId: 1
%% type: 单选题
%% chapter: 交流电
%% point: 变压器
%% method: 无
%% score: 3
%% degree: 550
%% duplicateId: 0
%% body:
如图所示，理想变压器的原、副线圈匝数比为 $2:1$，原线圈接在 $u=U_{\mathrm{m}}\sin(\omega t+\varphi)$ 的交流电源上，副线圈接定值电阻 $R_{0}$ 和滑动变阻器 $R$，电流表、电压表均为理想电表。滑动变阻器 $R$ 的滑片由 $a$ 端向 $b$ 端缓慢滑动 \xzanswer{C}
\onepicture[w=6.5cm]{figs/07.png}
\fourchoices[answer=C]
{电压表 $\mathrm{V}_{1}$、$\mathrm{V}_{2}$ 的示数之比为 $2:1$}
{电流表 $\mathrm{A}_{1}$ 示数不变}
{电流表 $\mathrm{A}_{2}$ 示数增加}
{变压器输出功率减小}

%% answer: C
%% memo:
\memoanswer{
A．理想变压器原副线圈的电压比等于匝数比，即原线圈两端电压（也就是 $\mathrm{V}_{1}$ 的示数）和副线圈两端总电压的比值为 $2:1$。但 $\mathrm{V}_{2}$ 测量的是滑动变阻器 $R$ 的分压，并非副线圈的总电压，因此 $\mathrm{V}_{1}$、$\mathrm{V}_{2}$ 的示数之比不等于 $2:1$，A 错误；

C．滑片从 $a$ 向 $b$ 滑动时，滑动变阻器接入电路的阻值减小。由于原线圈接的交流电源有效值恒定，根据变压规律 $\frac{U_{1}}{U_{\text{副总}}}=\frac{n_{1}}{n_{2}}$，副线圈的总电压有效值保持不变，副线圈回路总电阻为 $R_{0}+R$，总电阻减小，由欧姆定律可知副线圈总电流（即 $\mathrm{A}_{2}$ 的示数）增大，C 正确；

B．根据变流规律 $\frac{I_{1}}{I_{2}}=\frac{n_{2}}{n_{1}}$，副线圈电流 $I_{2}$ 增大时，原线圈电流 $I_{1}$（即 $\mathrm{A}_{1}$ 的示数）也会增大，B 错误；

D．变压器的输出功率 $P_{\text{出}}=U_{\text{副总}}I_{2}$。$U_{\text{副总}}$ 不变、$I_{2}$ 增大，因此输出功率增大，D 错误。

故选 C。
}
\item
%% number: 8
%% paperName: 2026年北京高考第 8 题 3 分
%% typeId: 1
%% type: 单选题
%% chapter: 曲线运动
%% point: 人造地球卫星
%% method: 无
%% score: 3
%% degree: 650
%% duplicateId: 0
%% body:
2026 年 2 月，我国捷龙三号运载火箭成功发射七颗卫星。假设某卫星在距离地面高度为 $h$ 的轨道上绕地球做匀速圆周运动，周期为 $T$。地球的半径为 $R$，下列说法正确的是 \xzanswer{A}
\fourchoices[answer=A]
{若已知引力常量 $G$，可计算地球的质量}
{该卫星的运行速度大于第一宇宙速度}
{该卫星的加速度大于地球表面的重力加速度}
{$h$ 越大，卫星运行的周期 $T$ 越小}

%% answer: A
%% memo:
\memoanswer{
A．卫星绕地球做匀速圆周运动，万有引力提供向心力，有
\[ G\frac{Mm}{(R+h)^{2}}=m\frac{4\pi^{2}}{T^{2}}(R+h) \]
约去卫星质量 $m$ 可得地球质量 $M=\frac{4\pi^{2}(R+h)^{3}}{GT^{2}}$，
已知 $G$、$T$、$R$、$h$ 即可计算地球质量，故 A 正确；

B．根据 $G\frac{Mm}{r^{2}}=m\frac{v^{2}}{r}$，可得卫星环绕速度公式为 $v=\sqrt{\frac{GM}{r}}$，
其中 $r$ 为轨道半径，第一宇宙速度是近地卫星（$r=R$）的环绕速度，是所有地球卫星的最大环绕速度，该卫星轨道半径 $r=R+h>R$，故运行速度小于第一宇宙速度，B 错误；

C．根据 $G\frac{Mm}{r^{2}}=mg$，卫星的加速度 $a=\frac{GM}{r^{2}}$，地球表面重力加速度 $g=\frac{GM}{R^{2}}$，
由于 $r=R+h>R$，故 $a<g$，C 错误；

D．由开普勒第三定律 $\frac{r^{3}}{T^{2}}=k$（$k$ 为常量），轨道半径 $r=R+h$ 越大，周期 $T$ 越大，故 $h$ 越大 $T$ 越大，D 错误。

故选 A。
}
\item
%% number: 9
%% paperName: 2026年北京高考第 9 题 3 分
%% typeId: 1
%% type: 单选题
%% chapter: 电场
%% point: 电容
%% method: 无
%% score: 3
%% degree: 550
%% duplicateId: 0
%% body:
将直流电源 $E$、电阻 $R$、平行板电容器 $C$ 与开关 $\mathrm{S}$ 连接成如图所示的电路。闭合开关 $\mathrm{S}$，电路达到稳定状态。下列说法正确的是 \xzanswer{B}
\onepicture[w=4.0cm]{figs/09.png}
\fourchoices[answer=B]
{此过程电源输出的能量等于电容器存储的电场能}
{若增大电容器两极板间距离，有从 $b$ 到 $a$ 的电流流过电阻}
{若断开开关，减小电容器两极板间距离，极板间电场强度增大}
{若断开开关，在电容器两极板间插入电介质后，两极板间电势差增大}

%% answer: B
%% memo:
\memoanswer{
A．闭合开关给电容器充电的过程中，有一部分能量在充电过程中通过电阻 $R$ 转化为焦耳热耗散，因此电源输出能量大于电容器存储的电场能，A 错误；

B．闭合开关稳定后，电容器两端电压等于电源电动势 $E$ 保持不变。若增大极板间距，由电容决定式 $C=\frac{\varepsilon_{r}S}{4\pi kd}$ 可知 $C$ 减小，根据 $Q=CU$，电容器带电量 $Q$ 减小，电容器放电，即存在从 $b$ 到 $a$ 的电流，B 正确；

C．断开开关后，电容器带电量 $Q$ 保持不变。极板间场强
\[ E_{\text{场}}=\frac{U}{d}=\frac{Q}{Cd}=\frac{Q}{\frac{\varepsilon_{r}S}{4\pi kd}\cdot d}=\frac{4\pi kQ}{\varepsilon_{r}S} \]
可见场强和极板间距 $d$ 无关，因此减小极板间距时电场强度不变，C 错误；

D．断开开关 $Q$ 不变，插入电介质后 $\varepsilon_{r}$ 增大，由 $C=\frac{\varepsilon_{r}S}{4\pi kd}$ 可知电容 $C$ 增大，根据 $U=\frac{Q}{C}$，两极板间电势差 $U$ 会减小，D 错误。

故选 B。
}
\item
%% number: 10
%% paperName: 2026年北京高考第 10 题 3 分
%% typeId: 1
%% type: 单选题
%% chapter: 运动和力
%% point: 变加速直线运动
%% method: 无
%% score: 3
%% degree: 450
%% duplicateId: 0
%% body:
将蹦极运动简化为直线运动。如图所示，轻质弹性绳的一端固定在 $O$ 点，另一端和人（视为质点）相连。人从 $a$ 点静止下落，经 $b$ 点时绳恰好伸直，到达最低点 $c$ 后向上运动。假设整个过程弹性绳的弹力与伸长量成正比，不计空气阻力。下列说法正确的是 \xzanswer{D}
\onepicture[h=4.8cm]{figs/10.png}
\fourchoices[answer=D]
{人运动到 $bc$ 中点处速度最大}
{人在 $b$ 点的加速度大于在 $c$ 点的加速度}
{从 $b$ 点到 $c$ 点的过程中，人一直处于失重状态}
{从 $a$ 点到 $c$ 点的过程中，重力势能与弹性势能之和先减小后增大}

%% answer: D
%% memo:
\memoanswer{
A．人速度最大的位置在合力为零处，即重力与弹力平衡的位置（设为 $d$ 点），此时 $mg=kx_{0}$（$x_{0}$ 为 $b$ 到 $d$ 的距离）。

若人从 $b$ 点静止释放，根据简谐运动对称性，最低点应在 $b$ 点下方 $2x_{0}$ 处；

现人从 $a$ 点下落，到达 $b$ 点已有速度，故最低点 $c$ 在 $b$ 点下方大于 $2x_{0}$ 处，即 $bc>2x_{0}$。

$bc$ 中点处的伸长量 $\Delta x=\frac{bc}{2}>x_{0}$，此时弹力 $F=k\Delta x>kx_{0}=mg$，合力向上，人处于减速阶段，故速度最大位置在 $bc$ 中点上方，故 A 错误；

B．在 $b$ 点，绳刚伸直，弹力为零，人只受重力，加速度 $a_{b}=g$。

在 $c$ 点，由 A 项分析可知伸长量 $bc>\frac{2mg}{k}$，弹力 $F_{c}=k\cdot bc>2mg$，

根据牛顿第二定律 $F_{c}-mg=ma_{c}$，

可得 $ma_{c}>mg$，即 $a_{c}>g$，所以 $c$ 点加速度大于 $b$ 点加速度，故 B 错误；

C．从 $b$ 点到 $c$ 点，人先做加速度减小的加速运动（$mg>F$，加速度向下，人处于失重状态），过平衡位置后做加速度增大的减速运动（$F>mg$，加速度向上，人处于超重状态），故人先处于失重状态后处于超重状态，故 C 错误；

D．对于人、地球和弹性绳组成的系统，只有重力和弹力做功，机械能守恒，从 $a$ 点到 $c$ 点的过程中，人的动能先增大后减小，所以重力势能与弹性势能之和先减小后增大，故 D 正确。

故选 D。
}
\item
%% number: 11
%% paperName: 2026年北京高考第 11 题 4 分
%% typeId: 1
%% type: 单选题
%% chapter: 磁场
%% point: 左手定则
%% method: 无
%% score: 4
%% degree: 450
%% duplicateId: 0
%% body:
如图所示实验装置中，两根平行直导线两端固定。通以电流，可演示电流间的相互作用，有助于建立电流产生磁场的认识。若可分别改变每根导线中的电流，还可以探究 \xzanswer{B}
\onepicture[h=4.8cm]{figs/11.png}
\fourchoices[answer=B]
{磁场强弱与距离的定性关系}
{磁场强弱与电流大小的定性关系}
{磁场是否可以叠加}
{磁场有无沿电流方向的分量}

%% answer: B
%% memo:
\memoanswer{
A．两根导线的间距不可改变，不可以定性得到磁场强弱和距离的关系，A 错误；

B．保持两根导线的间距不变，改变其中一根产生磁场的电流大小，通过另一根导线的受力形变程度变化，可以定性得到磁场强弱和电流大小的关系，B 正确；

C．该装置中只有两根平行通电直导线的磁场，没有额外的磁场叠加对照场景，无法验证磁场是否可以叠加，C 错误；

D．电流方向沿导线轴向，而安培力的方向始终垂直于电流方向，无法通过该装置检测是否存在沿电流方向的磁场分量，D 错误。

故选 B。
}
\item
%% number: 12
%% paperName: 2026年北京高考第 12 题 3 分
%% typeId: 1
%% type: 单选题
%% chapter: 电磁感应
%% point: 电磁感应中图象
%% method: 无
%% score: 3
%% degree: 450
%% duplicateId: 0
%% body:
用如图~\ref{2026北京12a} 所示的装置研究线圈中的感应电流。固定在支架上的线圈与电流传感器、滑动变阻器连接。$t=0$ 时磁铁从线圈上方某高处开始下落，穿过线圈并远离线圈。电流传感器采集的电流随时间变化的图像如图~\ref{2026北京12b} 所示。下列说法正确的是 \xzanswer{C}
\twopicture[numA=1, numB=2, labelA=2026北京12a, labelB=2026北京12b, h=3.0cm]
{figs/12a.png}
{figs/12b.png}
\fourchoices[answer=C]
{$t_{1}\sim t_{2}$ 时间内，感应电流方向从 $b$ 到 $a$}
{$t_{2}$ 时刻，穿过线圈的磁通量为 $0$}
{$t_{1}\sim t_{2}$、$t_{2}\sim t_{3}$ 两段时间内曲线与时间轴围成的面积相等}
{向左移动滑片 P，重复实验所得图像中的 $\Delta t=t_{2}-t_{1}$ 变小}

%% answer: C
%% memo:
\memoanswer{
A．$t_{1}\sim t_{2}$ 时间内，磁铁 N 极向下进入线圈，穿过线圈的磁通量向下增加。根据楞次定律，感应电流的磁场方向向上。由安培定则及图中线圈绕向可知，线圈中电流方向从上端流向下端，即外部电路中电流方向从 $a$ 到 $b$，故 A 错误；

B．$t_{2}$ 时刻，感应电流为 $0$，说明感应电动势为 $0$，即磁通量的变化率为 $0$。此时磁铁中心位于线圈中心，穿过线圈的磁通量最大，故 B 错误；

C．根据法拉第电磁感应定律和欧姆定律可得电荷量 $q=I\Delta t=n\frac{\Delta\Phi}{R}$。$t_{1}\sim t_{2}$ 时间内磁铁进入线圈，$t_{2}\sim t_{3}$ 时间内磁铁穿出线圈，两过程中磁通量变化量的大小 $\Delta\Phi$ 相等，故通过线圈的电荷量相等，即 $I-t$ 图像中曲线与时间轴围成的面积相等，故 C 正确；

D．向左移动滑片 P，滑动变阻器接入电路的阻值减小，回路总电阻 $R$ 减小。感应电流 $I$ 增大，线圈对磁铁的安培力阻碍作用增强，磁铁下落的加速度减小，平均速度减小，穿过线圈的时间 $\Delta t$ 变大，故 D 错误。

故选 C。
}
\item
%% number: 13
%% paperName: 2026年北京高考第 13 题 3 分
%% typeId: 1
%% type: 单选题
%% chapter: 光学
%% point: 光电效应
%% method: 无
%% score: 3
%% degree: 450
%% duplicateId: 0
%% body:
光照射到某半导体上，光子能量超过特定值 $W_{1}$，会激发半导体内部一些受束缚的电子，使其能够在半导体内部自由运动，从而增强导电能力，这属于光电导效应。可用如图所示电路研究光电导效应。用频率为 $\nu$ 的单色光分别均匀照射该半导体和逸出功为 $W_{2}$ 的金属，$W_{1}<h\nu<W_{2}$，$h$ 为普朗克常量。下列说法\textbf{错误}的是 \xzanswer{D}
\onepicture[w=4.7cm]{\includesvg{figs/13}}
\fourchoices[answer=D]
{金属不发生光电效应，半导体发生光电导效应}
{增大单色光的强度，电流表示数增大}
{对调电源的正负极，通过半导体的电流反向}
{半导体两端电压为 $0$ 时，电流表示数不为 $0$}

%% answer: D
%% memo:
\memoanswer{
A．光电效应的发生条件是光子能量 $h\nu>W_{2}$（金属逸出功），因此金属不会发生光电效应；而光电导效应的发生条件是光子能量 $h\nu>W_{1}$，因此光照射到半导体上时，会激发半导体内部受束缚的电子发生光电导效应。故 A 正确，不符合题意；

B．增大单色光强度，意味着单位时间内入射的光子数增多，半导体内部被激发的自由电子数量随之增加，半导体的电阻减小，则电路中电流会增大，所以电流表示数增大。故 B 正确，不符合题意；

C．当对调电源的正负极后，半导体两端的电场方向变成反向，半导体内部自由电子的受力方向也变成反向，则电子的定向移动方向随之反向，因此通过半导体的电流反向，故 C 正确，不符合题意；

D．该光电导效应仅在半导体内部激发自由载流子，不存在逸出半导体的光电子。当半导体两端电压为 $0$ 时，内部没有外电场驱动自由载流子定向移动，因此电流表示数为 $0$，故 D 错误，符合题意。

故选 D。
}
\item
%% number: 14
%% paperName: 2026年北京高考第 14 题 3 分
%% typeId: 1
%% type: 单选题
%% chapter: 其他
%% point: 物理思想
%% method: 类比
%% score: 3
%% degree: 450
%% duplicateId: 0
%% body:
霍普菲尔德模型是神经网络实现联想记忆的核心机制之一。作为物理思想方法在人工神经网络领域应用的重要成果，这项工作获得 2024 年诺贝尔物理学奖。

简单来说，该模型以 $0$ 或 $1$ 表示神经元的状态，所有神经元状态的集合表示神经网络状态；并给神经网络定义了一个“能量函数”（以下简称函数），通过训练使函数变化产生多个极小值，每个极小值对应的神经网络状态代表一个存储信息，也就是记忆。在信息检索与重现过程中，神经网络状态向函数值降低的方向演化，最终达到某个极小值为止，实现已存储信息的重现。该过程可做如下类比：在如图所示的峰谷交错地形中某位置释放一个小球，小球始终沿着重力势能减小的方向运动，直至停在某个山谷底部；小球的位置对应神经网络状态，势能对应函数值。例如，输入模糊字母对应设定小球的初始位置 $a$，小球到达谷底 $b$ 对应重现出已存储的字母 J。下列说法正确的是 \xzanswer{B}
\onepicture[w=9.5cm]{\includesvg{figs/14}}
\fourchoices[answer=B]
{信息的重现过程可类比小球在光滑的凹凸曲面上运动}
{输入具有微小差别的两个模糊字母，可能重现出不同的已存储字母}
{输入函数值相等的模糊字母，都可以重现出相同的已存储字母}
{函数变化使某个谷底消失，仍可重现出原来对应的存储信息}

%% answer: B
%% memo:
\memoanswer{
A．若曲面光滑，小球运动过程中动能和重力势能相互转化，小球会在山谷间往复运动，无法最终停在谷底。而信息重现过程要求小球最终停在势能极小值点，类比的是存在阻力、小球沿势能持续减小方向运动直到静止的过程，A 错误；

B．若两个有微小差别的初始位置（对应微小差别的模糊字母），恰好分属两个不同山谷的势力范围，小球会沿各自的路径滑入不同的谷底，对应重现出不同的已存储字母，因此该情况是可能存在的，B 正确；

C．函数值相等仅对应初始位置的重力势能相等，两个势能相等的初始位置完全可能分属不同山谷的吸引域，最终滑入不同谷底，重现出不同的存储信息，C 错误；

D．每个谷底对应一个已存储的记忆信息，若该谷底消失，不存在对应的势能极小值点，无法重现原来对应的存储信息，D 错误。

故选 B。
}
\item
%% number: 15
%% paperName: 2026年北京高考第 15 题 8 分
%% typeId: 4
%% type: 实验题
%% chapter: 运动和力
%% point: 验证牛顿第二定律
%% method: 无
%% score: 8
%% degree: 650
%% duplicateId: 0
%% body:
\begin{enumerate}
	\item
	图~\ref{2026北京15a} 为“探究加速度与力的关系”的实验装置示意图。平衡阻力后，认为桶和砂所受的重力等于小车所受的合力，则桶和砂的总质量 $m$ 与小车的质量 $M$ 满足的关系是 \tkanswer{$m\ll M$}。
	\item
	下列实验操作，正确的是 \tkanswer{AC}（填选项前的字母）。
	\threechoices[answer=AC]
	{在“验证机械能守恒定律”实验中，先接通电源，后释放重锤}
	{在“探究气体等温变化规律”实验中，快速拉动柱塞改变空气柱体积}
	{在“测量玻璃的折射率”实验中，为准确确定入射光线的方向，两枚大头针的距离应适当大些}
	\item
	在“测量金属丝的电阻率”实验中，某同学将实验器材连接成如图~\ref{2026北京15b} 所示的电路。闭合开关后，电流表和电压表均有示数，但移动滑动变阻器的滑片，电流表和电压表的示数几乎不变。拆除导线 \tkanswer{$c$ 或 $ac$} 可排除故障（填图中导线字母编号）。

	排除故障后进行实验。仅考虑由电表内阻引起的实验误差，电阻率的测量值 \tkanswer{偏小}（选填“偏大”、“偏小”或“没有影响”）。
\end{enumerate}
\twopicture[numA=1, numB=2, labelA=2026北京15a, labelB=2026北京15b, h=3.7cm, align=right]
{figs/15a.png}
{figs/15b.png}
%% answer:
\jdanswer{
\begin{enumerate}
	\item
	$m\ll M$

	\item
	AC

	\item
	$c$ 或 $ac$，偏小
\end{enumerate}
}
%% memo:
\memoanswer{
\begin{enumerate}
	\item
	对小车和桶砂整体由牛顿第二定律得加速度 $a=\frac{mg}{M+m}$，小车实际受到的拉力
	\[ T=Ma=\frac{Mmg}{M+m}=\frac{mg}{1+\frac{m}{M}} \]
	只有当 $m\ll M$ 时，才可近似认为拉力 $T$ 等于桶和砂的总重力 $mg$。

	\item
	A．使用打点计时器的实验都需遵循“先接通电源，待打点稳定后再释放重物”的操作规范以保证充分利用纸带，A 正确；

	B．探究气体等温变化规律时，需要保证气体状态变化过程温度始终和环境温度一致，快速拉动柱塞会使气体温度发生变化，违背等温条件，B 错误；

	C．测量玻璃折射率时，两枚大头针间距适当大些，可以减小确定入射光线方向的角度误差，提高实验精度，C 正确。

	故选 AC。

	\item
	故障分析：图中导线 $c$ 直接将滑动变阻器的右上接线柱和开关相连，此时滑动变阻器的全部电阻被接入电路，相当于定值电阻，因此移动滑片时电路总电阻几乎不变，电表示数几乎不变。拆除多余的导线 $c$，滑动变阻器即可恢复正常的调节功能，排除故障，也可以拆除导线 $ac$，故答案为 $c$ 或 $ac$。

	误差分析：该电路为电流表外接法，电压表测量的是电阻丝的真实电压，但电流表测量的是电阻丝和电压表的总电流，电流测量值偏大。根据 $R=\frac{U}{I}$，电阻测量值偏小，再由电阻率公式 $\rho=\frac{RS}{L}$，因此电阻率的测量值偏小。
\end{enumerate}
}
\item
%% number: 16
%% paperName: 2026年北京高考第 16 题 10 分
%% typeId: 4
%% type: 实验题
%% chapter: 曲线运动
%% point: 研究平抛运动实验
%% method: 无
%% score: 10
%% degree: 450
%% duplicateId: 0
%% body:
用如图~\ref{2026北京16a} 所示装置研究平抛运动。将白纸和复写纸固定在背板上。钢球沿斜槽轨道滑下后飞出，落在水平挡板上，钢球侧面会在白纸上挤压出一个印迹。移动挡板，重复实验，白纸上留下一系列印迹。
\begin{enumerate}
	\item
	下列实验操作必要的是 \tkanswer{AB}（填选项前的字母）。
	\threechoices[answer=AB]
	{背板调节到竖直平面内}
	{斜槽轨道的末端保持水平}
	{等间距移动挡板}
	\item
	实验操作时，钢球每次从同一个位置由静止释放，目的是 \tkanswer{使钢球每次从斜槽轨道的末端飞出的速度相同}。
	\item
	用平滑曲线将白纸上的印迹连接起来，如图~\ref{2026北京16b} 所示。以抛出点为坐标原点 $O$，以竖直方向和水平方向分别作为 $y$ 轴和 $x$ 轴，建立 $Oxy$ 坐标系。已知平抛运动竖直方向分运动为自由落体运动，为探究水平分运动的规律，需要确定“相等的时间间隔”。简要说明如何在轨迹上选点来确定相等的时间间隔。
	\item
	某同学实验时忘了标记抛出点的位置。为解决该问题，他以轨迹上某点为坐标原点 $O'$，建立 $O'xy$ 坐标系。在轨迹上取几个点，测量各点的坐标（$x$，$y$）。作出的 $\frac{y}{x}-x$ 图像为一条直线，斜率为 $k$，截距为 $b$，如图~\ref{2026北京16c} 所示。根据平抛运动规律，可得在 $O'xy$ 坐标系中钢球做平抛运动的抛出点坐标为 \tkanswer{$\left(-\frac{b}{2k},-\frac{b^{2}}{4k}\right)$}。
\end{enumerate}
\threepicture[numA=1, numB=2, numC=3, labelA=2026北京16a, labelB=2026北京16b, labelC=2026北京16c, h=4.0cm, align=right]
{figs/16a.png}
{figs/16b.png}
{figs/16c.png}
%% answer:
\jdanswer{
\begin{enumerate}
	\item
	AB

	\item
	使钢球每次从斜槽轨道的末端飞出的速度相同

	\item
	因为平抛竖直分运动是初速度为 $0$ 的自由落体运动，满足连续相等时间内的竖直位移之比为 $1:3:5:\cdots$（或位移差恒定）；从抛出点 $O$ 开始，在 $y$ 轴上取纵坐标分别为 $y$、$4y$、$9y$、$\cdots$ 的点（即竖直位移满足 $y_{1}:y_{2}:y_{3}:\cdots=1:4:9:\cdots$），过这些点作平行于 $x$ 轴的水平线，与轨迹的交点对应的时间间隔就是相等的时间间隔。（或：在轨迹上取点，使相邻两点的竖直方向位移差恒定，对应的时间间隔也相等。）

	\item
	$\left(-\frac{b}{2k},-\frac{b^{2}}{4k}\right)$
\end{enumerate}
}
%% memo:
\memoanswer{
\begin{enumerate}
	\item
	A．钢球的平抛运动轨迹在竖直平面内，为了印迹能准确记录运动位置，应将背板调节到竖直平面内，故 A 正确；

	B．为了保证小球抛出时的速度处于水平方向，应调节斜槽轨道的末端保持水平，故 B 正确；

	C．挡板不需要等间距移动，只要能在轨迹上得到多个分布合理的点即可，故 C 错误。故选 AB。

	\item
	实验操作时，钢球每次从同一个位置由静止释放，目的是：使钢球每次从斜槽轨道的末端飞出的速度相同。

	\item
	因为平抛竖直分运动是初速度为 $0$ 的自由落体运动，满足连续相等时间内的竖直位移之比为 $1:3:5:\cdots$（或位移差恒定）；从抛出点 $O$ 开始，在 $y$ 轴上取纵坐标分别为 $y$、$4y$、$9y$、$\cdots$ 的点（即竖直位移满足 $y_{1}:y_{2}:y_{3}:\cdots=1:4:9:\cdots$），过这些点作平行于 $x$ 轴的水平线，与轨迹的交点对应的时间间隔就是相等的时间间隔。（或：在轨迹上取点，使相邻两点的竖直方向位移差恒定，对应的时间间隔也相等。）

	\item
	设钢球在 $O'$ 坐标系中，经过坐标（$x$，$y$）的时刻为 $t$，平抛初速度为 $v_{0}$，抛出点在 $O'$ 系中的坐标为（$x_{0}$，$y_{0}$）。根据平抛运动规律有 $x-x_{0}=v_{0}t$，$y-y_{0}=\frac{1}{2}gt^{2}$。

	联立整理可得 $y=\frac{g}{2v_{0}^{2}}x^{2}-\frac{gx_{0}}{v_{0}^{2}}x+\frac{gx_{0}^{2}}{2v_{0}^{2}}+y_{0}$。

	已知作出的 $\frac{y}{x}-x$ 图像为一条直线，斜率为 $k$，截距为 $b$，则有 $\frac{y}{x}=kx+b$，整理可得 $y=kx^{2}+bx$。

	联立可得 $\frac{g}{2v_{0}^{2}}=k$，$-\frac{gx_{0}}{v_{0}^{2}}=b$，$\frac{gx_{0}^{2}}{2v_{0}^{2}}+y_{0}=0$。

	解得 $x_{0}=-\frac{b}{2k}$，$y_{0}=-\frac{b^{2}}{4k}$。

	则在 $O'xy$ 坐标系中钢球做平抛运动的抛出点坐标为 $\left(-\frac{b}{2k},-\frac{b^{2}}{4k}\right)$。
\end{enumerate}
}
\item
%% number: 17
%% paperName: 2026年北京高考第 17 题 9 分
%% typeId: 6
%% type: 计算题
%% chapter: 电磁感应
%% point: 电磁感应中的变加速
%% method: 无
%% score: 9
%% degree: 650
%% duplicateId: 0
%% body:
如图所示，匀强磁场的磁感应强度大小为 $B$、方向竖直向下，宽为 $L$ 的 U 型长直金属导轨水平放置，导轨左端接阻值为 $R$ 的电阻。质量为 $m$ 的导体棒 $ab$ 在导轨上以初速度 $v_{0}$ 开始运动。运动过程中导体棒与导轨垂直且接触良好，不计导体棒和导轨的电阻及它们间的摩擦。求：
\begin{enumerate}
	\item
	初始时刻通过导体棒的电流大小和方向；
	\item
	初始时刻导体棒加速度的大小；
	\item
	导体棒速度从 $v_{0}$ 减为 $0$ 的过程中电阻 $R$ 产生的热量。
\end{enumerate}
\onepicture[w=6.2cm, align=right]{\includesvg{figs/17}}
%% answer:
\jdanswer{
\begin{enumerate}
	\item
	导体棒的电流方向 $a\to b$，大小为 $\frac{BLv_{0}}{R}$

	\item
	$\frac{B^{2}L^{2}v_{0}}{mR}$

	\item
	$\frac{1}{2}mv_{0}^{2}$
\end{enumerate}
}
%% memo:
\memoanswer{
（1）根据右手定则可知，导体棒中感应电流的方向为 $a\to b$；刚开始运动时，感应电动势为 $E=BLv_{0}$，

通过导体棒的电流大小 $I=\frac{E}{R}=\frac{BLv_{0}}{R}$。

（2）初始时刻导体棒所受安培力为 $F=BIL=\frac{B^{2}L^{2}v_{0}}{R}$，

由牛顿第二定律有 $F=ma$，

解得 $a=\frac{B^{2}L^{2}v_{0}}{mR}$。

（3）根据题意，由能量守恒定律可知，电阻 $R$ 产生的热量 $Q=\frac{1}{2}mv_{0}^{2}$。
}
\item
%% number: 18
%% paperName: 2026年北京高考第 18 题 9 分
%% typeId: 6
%% type: 计算题
%% chapter: 机械振动
%% point: 单摆
%% method: 无
%% score: 9
%% degree: 550
%% duplicateId: 0
%% body:
长为 $l$ 的轻绳一端固定在 $O$ 点，另一端系一个质量为 $m$ 的小球。重力加速度为 $g$，不计空气阻力。
\begin{enumerate}
	\item
	如图~\ref{2026北京18a} 所示，将小球拉至某一位置，轻绳与竖直方向夹角为 $\theta$，静止释放小球。
	a．求小球运动到最低点时受到绳的拉力大小 $F$；
	b．若 $\theta$ 较小，释放后小球做简谐运动。求一个周期内绳对小球拉力的冲量大小 $I$。
	\item
	如图~\ref{2026北京18b} 所示，使小球在水平面内做匀速圆周运动，轻绳与竖直方向夹角为 $\theta$，求小球运动的周期，并与（1）b 中小球运动的周期比较大小。
\end{enumerate}
\twopicture[numA=1, numB=2, labelA=2026北京18a, labelB=2026北京18b, h=5.0cm, align=right]
{figs/18a.png}
{figs/18b.png}
%% answer:
\jdanswer{
\begin{enumerate}
	\item
	a．$F=mg(3-2\cos\theta)$；b．$I=2\pi m\sqrt{gl}$

	\item
	$T'=2\pi\sqrt{\frac{l\cos\theta}{g}}$，小于（1）b 中小球运动的周期
\end{enumerate}
}
%% memo:
\memoanswer{
（1）a．小球从静止释放到最低点的过程中，根据机械能守恒可得
\[ mgl(1-\cos\theta)=\frac{1}{2}mv^{2} \]
解得 $v=\sqrt{2gl(1-\cos\theta)}$。

在最低点，由牛顿第二定律可得 $F-mg=m\frac{v^{2}}{l}$，

解得 $F=mg(3-2\cos\theta)$。

b．若 $\theta$ 较小，释放后小球做简谐运动，根据单摆周期公式可得 $T=2\pi\sqrt{\frac{l}{g}}$。

根据动量定理可得 $I-mgT=\Delta p=0$，

联立解得一个周期内绳对小球拉力的冲量大小为 $I=2\pi m\sqrt{gl}$。

（2）使小球在水平面内做匀速圆周运动，轻绳与竖直方向夹角为 $\theta$，根据牛顿第二定律可得 $mg\tan\theta=m\frac{4\pi^{2}}{T^{2}}r$，

其中 $r=l\sin\theta$，

解得小球运动的周期为 $T'=2\pi\sqrt{\frac{l\cos\theta}{g}}$。

而（1）b 中小球运动的周期为 $T=2\pi\sqrt{\frac{l}{g}}$，

由于 $\cos\theta<1$，可得 $T'<T$，即（2）中小球运动的周期小于（1）b 中小球运动的周期。
}
\item
%% number: 19
%% paperName: 2026年北京高考第 19 题 10 分
%% typeId: 6
%% type: 计算题
%% chapter: 磁场
%% point: 质谱仪回旋加速器
%% method: 无
%% score: 10
%% degree: 450
%% duplicateId: 0
%% body:
药物研发中常需要二级质谱分析。一级质谱分析只能检出与目标离子具有相同比荷的离子（母离子）。为进一步鉴别目标离子，需将母离子打碎成更小的离子（子离子），进行第二次质谱分析。某种二级串联质谱系统装置如图所示，假设离子均带正电。
\begin{enumerate}
	\item
	第一级质谱：A 区中的离子由静止经电场加速后，垂直进入磁感应强度为 $B$ 的匀强磁场，比荷为 $\frac{q}{m}$ 的母离子做半径为 $R$ 的匀速圆周运动进入 C 区。求加速电压 $U_{1}$。
	\item
	两级质谱连接：C 区产生的子离子经聚焦、冷却后，由静止经电压 $U_{2}$ 加速后进入 D 区。为保证不同比荷的子离子均由小孔 S 垂直进入第二级质谱的磁场，在 D 区施加某种场，使这些子离子都做以 $O$ 点为圆心、半径为 $r$ 的匀速圆周运动。请设计 D 区中的场：
	a．推导说明仅施加匀强磁场是否可行；
	b．推导说明仅施加电场是否可行。
\end{enumerate}
\onepicture[w=8.8cm, align=right]{figs/19.png}
%% answer:
\jdanswer{
\begin{enumerate}
	\item
	$U_{1}=\frac{qB^{2}R^{2}}{2m}$

	\item
	a．仅施加匀强磁场不可行。子离子经电压加速，由动能定理得 $qU_{2}=\frac{1}{2}mv'^{2}$；若仅加匀强磁场，洛伦兹力提供向心力，$qv'B=m\frac{v'^{2}}{r}$，得 $r=\frac{mv'}{qB}=\frac{1}{B}\sqrt{\frac{2mU_{2}}{q}}$，可见轨道半径 $r$ 与子离子的比荷 $\frac{q}{m}$ 有关，无法保证所有子离子都以固定半径 $r$ 做圆周运动，因此仅施加匀强磁场不可行。

	b．仅施加电场可行，需施加以 $O$ 为圆心的径向辐向电场。由 $qE=m\frac{v'^{2}}{r}$ 与 $qU_{2}=\frac{1}{2}mv'^{2}$ 可得 $E=\frac{mv'^{2}}{qr}=\frac{2qU_{2}}{qr}=\frac{2U_{2}}{r}$，该结果与比荷无关，故可行。
\end{enumerate}
}
%% memo:
\memoanswer{
（1）母离子在 A 区由静止经电场加速，由动能定理得 $qU_{1}=\frac{1}{2}mv^{2}$。

母离子进入匀强磁场后，洛伦兹力提供匀速圆周运动的向心力 $qvB=m\frac{v^{2}}{R}$，

联立可得 $U_{1}=\frac{qB^{2}R^{2}}{2m}$。

（2）a．仅施加匀强磁场不可行。

子离子经电压加速，由动能定理得 $qU_{2}=\frac{1}{2}mv'^{2}$。

若仅加匀强磁场，洛伦兹力提供向心力，轨道半径满足 $qv'B=m\frac{v'^{2}}{r}$，

得 $r=\frac{mv'}{qB}=\frac{1}{B}\sqrt{\frac{2mU_{2}}{q}}$。

可见轨道半径 $r$ 与子离子的比荷 $\frac{q}{m}$ 有关，不同比荷的子离子加速后速度不同，在同一匀强磁场中做圆周运动的半径各不相同，无法保证所有子离子都以固定半径 $r$ 做圆周运动，因此仅施加匀强磁场不可行。

b．仅施加电场可行，需施加以 $O$ 为圆心的径向辐向电场。

若仅施加电场，要让所有子离子都以 $O$ 为圆心、$r$ 为半径做匀速圆周运动，电场力需完全提供向心力，且电场方向始终指向圆心。

设距离 $O$ 点 $r$ 处的电场强度大小为 $E$，对任意子离子有 $qE=m\frac{v'^{2}}{r}$，

结合子离子的加速动能关系 $qU_{2}=\frac{1}{2}mv'^{2}$，

代入得 $E=\frac{mv'^{2}}{qr}=\frac{2qU_{2}}{qr}=\frac{2U_{2}}{r}$。

该结果与子离子的比荷 $\frac{q}{m}$ 无关，只要 D 区是辐向电场，在半径为 $r$ 的圆周上各处电场强度大小恒为 $\frac{2U_{2}}{r}$、方向沿径向指向圆心，不同比荷的子离子就都能获得大小合适的向心力，完成半径为 $r$ 的匀速圆周运动，最终从 S 垂直射出，因此仅施加该种径向电场是可行的。
}
\item
%% number: 20
%% paperName: 2026年北京高考第 20 题 12 分
%% typeId: 6
%% type: 计算题
%% chapter: 运动和力
%% point: 动量和能量
%% method: 无
%% score: 12
%% degree: 350
%% duplicateId: 0
%% body:
中子或暗物质粒子等中性粒子因为不带电，无法直接通过电磁效应探测。
\begin{enumerate}
	\item
	实验上，通常让中性粒子与原子核碰撞，通过原子核获得的动能来探测中性粒子的性质。质量为 $m$、速度为 $v$ 的中性粒子与质量为 $M$ 的静止原子核发生对心碰撞。若碰撞后中性粒子与原子核的速度相同，求原子核动能 $E_{k}$。
	\item
	理论上，通过米格达尔效应也可以探测中性粒子。米格达尔效应可近似为：中性粒子与原子核碰撞，系统损失的动能都传递给原子核外电子，成为可探测的电子信号。
	a．质量为 $m$、速度为 $v$ 的中性粒子与质量为 $M$ 的静止原子核对心碰撞时发生了米格达尔效应，求电子可获得的最大能量 $E_{\text{电}}$；
	b．2026 年 1 月，我国科学家团队首次在实验中直接观测到中子与原子核碰撞过程中的米格达尔效应，这一发现为轻暗物质探测突破阈值瓶颈提供了关键支撑。若暗物质粒子的质量 $m=0.5\UGeVcSq$（$1\UGeVcSq\approx1.78\times10^{-27}\Ukg$），原子核的质量 $M=120\UGeVcSq$，暗物质粒子与原子核的各类碰撞中，原子核可获得的最大动能为 $0.03\UkeV$，探测器能分辨的最小能量值为 $1\UkeV$，请结合数据分析说明：与通过原子核获得的动能来探测中性粒子相比，米格达尔效应有助于探测质量更小的暗物质粒子。
\end{enumerate}
%% answer:
\jdanswer{
\begin{enumerate}
	\item
	$E_{k}=\frac{Mm^{2}v^{2}}{2(M+m)^{2}}$

	\item
	a．$E_{\text{电}}=\frac{Mmv^{2}}{2(M+m)}$；b．$E_{\text{电}}\approx1.8\UkeV$，大于 $1\UkeV$ 的探测阈值，故米格达尔效应有助于探测质量更小的暗物质粒子。
\end{enumerate}
}
%% memo:
\memoanswer{
（1）碰撞后两者速度相同，系统动量守恒 $mv=(M+m)v_{\text{共}}$，

解得共同速度 $v_{\text{共}}=\frac{mv}{M+m}$。

原子核动能 $E_{k}=\frac{1}{2}Mv_{\text{共}}^{2}=\frac{Mm^{2}v^{2}}{2(M+m)^{2}}$。

（2）a．中性粒子和原子核碰撞时，系统损失的动能全部转移给电子，因此电子获得的能量等于碰撞过程的动能损失。

对心碰撞中完全非弹性碰撞的动能损失最大，对应电子获得最大能量：初始总动能为 $E_{0}=\frac{1}{2}mv^{2}$，

完全非弹性碰撞后总动能为 $\frac{1}{2}(M+m)v_{\text{共}}^{2}$，

因此
\[ E_{\text{电}}=E_{0}-\frac{1}{2}(M+m)v_{\text{共}}^{2}=\frac{1}{2}mv^{2}-\frac{1}{2}(M+m)\left(\frac{mv}{M+m}\right)^{2}=\frac{Mmv^{2}}{2(M+m)} \]

b．原子核可获得的最大动能为 $0.03\UkeV$，远小于探测器最小分辨能量 $1\UkeV$，无法通过原子核动能直接探测。暗物质粒子与原子核的各类碰撞中，弹性碰撞时，原子核获得的动能最大，根据动量守恒和能量守恒有 $mv=mv_{1}+Mv_{2}$，$\frac{1}{2}mv^{2}=\frac{1}{2}mv_{1}^{2}+\frac{1}{2}Mv_{2}^{2}$，

可得 $E_{\text{kmax}}'=\frac{1}{2}M\left(\frac{2m}{M+m}v\right)^{2}=0.03\UkeV$ ①。

由 a 分析可知米格达尔效应中电子可获得的最大能量 $E_{\text{电}}=\frac{Mmv^{2}}{2(M+m)}$ ②。

联立①②可得 $E_{\text{电}}\approx1.8\UkeV$，该值大于 $1\UkeV$ 的探测阈值，可以被探测器捕捉。说明对于质量更小的暗物质粒子，直接碰撞转移给原子核的动能远低于探测阈值无法被识别，而米格达尔效应能将碰撞损失的动能集中转移给电子，使可探测的信号能量超过阈值，因此可以探测到质量更小的暗物质粒子。
}
\end{enumerate}

\end{document}
